\documentclass[12pt]{article}
\usepackage{amsmath,amssymb}

\begin{document}
\section*{{2023 AIME II Problems/Problem 11}}
Source: AoPS Wiki

\subsection*{Solution}
Denote by $\mathcal C$ a collection of 16 distinct subsets of $\left\{ 1, 2, 3, 4, 5 \right\}$ .
Denote $N = \min \left\{ |S|: S \in \mathcal C \right\}$ . Case 1: $N = 0$ . This entails $\emptyset \in \mathcal C$ .
Hence, for any other set $A \in \mathcal C$ , we have $\emptyset \cap A = \emptyset$ . This is infeasible. Case 2: $N = 1$ . Let $\{a_1\} \in \mathcal C$ .
To get $\{a_1\} \cap A \neq \emptyset$ for all $A \in \mathcal C$ .
We must have $a_1 \in \mathcal A$ . The total number of subsets of $\left\{ 1, 2, 3, 4, 5 \right\}$ that contain $a_1$ is $2^4 = 16$ .
Because $\mathcal C$ contains 16 subsets.
We must have $\mathcal C = \left\{ \{a_1\} \cup A : \forall \ A \subseteq \left\{ 1, 2, 3, 4, 5 \right\} \backslash \left\{a_1 \right\} \right\}$ .
Therefore, for any $X, Y \in \mathcal C$ , we must have $X \cap Y \supseteq \{a_1\}$ .
So this is feasible. Now, we count the number of $\mathcal C$ in this case.
We only need to determine $a_1$ .
Therefore, the number of solutions is 5. Case 3: $N = 2$ . Case 3.1: There is exactly one subset in $\mathcal C$ that contains 2 elements. Denote this subset as $\left\{ a_1, a_2 \right\}$ .
We then put all subsets of $\left\{ 1, 2, 3, 4, 5 \right\}$ that contain at least three elements into $\mathcal C$ , except $\left\{ a_3, a_4, a_5 \right\}$ .
This satisfies $X \cap Y \neq \emptyset$ for any $X, Y \in \mathcal C$ . Now, we count the number of $\mathcal C$ in this case.
We only need to determine $\left\{ a_1, a_2 \right\}$ .
Therefore, the number of solutions is $\binom{5}{2} = 10$ . Case 3.2: There are exactly two subsets in $\mathcal C$ that contain 2 elements. They must take the form $\left\{ a_1, a_2 \right\}$ and $\left\{ a_1, a_3 \right\}$ . We then put all subsets of $\left\{ 1, 2, 3, 4, 5 \right\}$ that contain at least three elements into $\mathcal C$ , except $\left\{ a_3, a_4, a_5 \right\}$ and $\left\{ a_2, a_4, a_5 \right\}$ .
This satisfies $X \cap Y \neq \emptyset$ for any $X, Y \in \mathcal C$ . Now, we count the number of $\mathcal C$ in this case.
We only need to determine $\left\{ a_1, a_2 \right\}$ and $\left\{ a_1, a_3 \right\}$ .
Therefore, the number of solutions is $5 \cdot \binom{4}{2} = 30$ . Case 3.3: There are exactly three subsets in $\mathcal C$ that contain 2 elements.
They take the form $\left\{ a_1, a_2 \right\}$ , $\left\{ a_1, a_3 \right\}$ , $\left\{ a_1, a_4 \right\}$ . We then put all subsets of $\left\{ 1, 2, 3, 4, 5 \right\}$ that contain at least three elements into $\mathcal C$ , except $\left\{ a_3, a_4, a_5 \right\}$ , $\left\{ a_2, a_4, a_5 \right\}$ , $\left\{ a_2, a_3, a_5 \right\}$ .
This satisfies $X \cap Y \neq \emptyset$ for any $X, Y \in \mathcal C$ . Now, we count the number of $\mathcal C$ in this case.
We only need to determine $\left\{ a_1, a_2 \right\}$ , $\left\{ a_1, a_3 \right\}$ , $\left\{ a_1, a_4 \right\}$ .
Therefore, the number of solutions is $5 \cdot \binom{4}{3} = 20$ . Case 3.4: There are exactly three subsets in $\mathcal C$ that contain 2 elements.
They take the form $\left\{ a_1, a_2 \right\}$ , $\left\{ a_1, a_3 \right\}$ , $\left\{ a_2, a_3 \right\}$ . We then put all subsets of $\left\{ 1, 2, 3, 4, 5 \right\}$ that contain at least three elements into $\mathcal C$ , except $\left\{ a_3, a_4, a_5 \right\}$ , $\left\{ a_2, a_4, a_5 \right\}$ , $\left\{ a_1, a_4, a_5 \right\}$ .
This satisfies $X \cap Y \neq \emptyset$ for any $X, Y \in \mathcal C$ . Now, we count the number of $\mathcal C$ in this case.
We only need to determine $\left\{ a_1, a_2 \right\}$ , $\left\{ a_1, a_3 \right\}$ , $\left\{ a_2, a_3 \right\}$ .
Therefore, the number of solutions is $\binom{5}{3} = 10$ . Case 3.5: There are exactly four subsets in $\mathcal C$ that contain 2 elements.
They take the form $\left\{ a_1, a_2 \right\}$ , $\left\{ a_1, a_3 \right\}$ , $\left\{ a_1, a_4 \right\}$ , $\left\{ a_1, a_5 \right\}$ . We then put all subsets of $\left\{ 1, 2, 3, 4, 5 \right\}$ that contain at least three elements into $\mathcal C$ , except $\left\{ a_3, a_4, a_5 \right\}$ , $\left\{ a_2, a_4, a_5 \right\}$ , $\left\{ a_1, a_4, a_5 \right\}$ , $\left\{ a_2, a_3, a_4 \right\}$ .
This satisfies $X \cap Y \neq \emptyset$ for any $X, Y \in \mathcal C$ . Now, we count the number of $\mathcal C$ in this case.
We only need to determine $\left\{ a_1, a_2 \right\}$ , $\left\{ a_1, a_3 \right\}$ , $\left\{ a_1, a_4 \right\}$ , $\left\{ a_1, a_5 \right\}$ .
Therefore, the number of solutions is 5. Putting all subcases together, the number of solutions is this case is $10 + 30 + 20 + 10 + 5 = 75$ . Case 4: $N \geq 3$ . The number of subsets of $\left\{ 1, 2, 3, 4, 5 \right\}$ that contain at least three elements is $\sum_{i=3}^5 \binom{5}{i} = 16$ .
Because $\mathcal C$ has 16 elements, we must select all such subsets into $\mathcal C$ .
Therefore, the number of solutions in this case is 1. Putting all cases together, the total number of $\mathcal C$ is $5 + 75 + 1 = \boxed{\textbf{(081) }}$ .

\end{document}
